\subsubsection{Input systém a zpracování uživatelského vstupu}

Vstupy od uživatele se zpracovávají pomocí Unity Input Systemu, který poskytuje modernější a flexibilnější přístup ke zpracování vstupů oproti staršímu Input Manageru. Pomocí atributu \texttt{[SerializeField] private InputActionReference} se přiřazují jednotlivé akce v editoru. Každá akce repreguje konkrétní typ vstupu, jako je pohyb, skok nebo střelba.

K registraci inputu dochází v metodě \texttt{OnEnable()}, kde se připojují callbacky na jednotlivé akce. Tyto callbacky jsou metody, které se automaticky volají při změně stavu vstupu. Příkladem je \texttt{OnMoveInput()} pro pohyb, \texttt{OnJumpInput()} pro skok a podobně. Při uvolnění skriptu nebo deaktivaci objektu se v \texttt{OnDisable()} tyto callbacky odpojují, což zabraňuje chybám při volání metod na neexistující objekty.

Tento přístup k input systému je výhodný, protože umožňuje snadnou rekonfiguraci vstupů bez nutnosti měnit kód. V editoru lze přiřadit libovolnou klávesu nebo tlačítko k akci a systém automaticky zpracuje změnu.

\begin{lstlisting}[language=csharp,caption={Listing 3: Deklarace vstupních akcí}]
[Header("input")]
public Vector2 moveInput { get; private set; }
[SerializeField] private InputActionReference moveAction;
[SerializeField] private InputActionReference jumpAction;
[SerializeField] private InputActionReference dashAction;
[SerializeField] private InputActionReference shootAction;

private void OnEnable()
{
    moveAction.action.Enable();
    jumpAction.action.Enable();
    shootAction.action.Enable();
    dashAction.action.Enable();

    jumpAction.action.performed += OnJumpInput;
    jumpAction.action.canceled += OnJumpRelease;
    dashAction.action.performed += OnDashInput;
    shootAction.action.performed += OnShootInput;
}
\end{lstlisting}
